% FUNDAMENTAÇÃO — Volume 2 (gerado por expandir_volume.py)

\chapter{O Método Híbrido Fônico-Kumon no 2º Ano}

No 2º ano, o método híbrido fônico-kumon avança da decodificação para a regulação ortográfica e a fluência. O treino diário e repetitivo (espírito Kumon) organiza o automatismo; a técnica fônica garante a correspondência grafema-fonema; a Técnica da Boquinha ancora a articulação correta de cada som. O professor conduz a sequência didática granular, e os pontos de rastreio integrado (SPEC-935-R211) entre unidades orientam a observação pedagógica sem diagnosticar.


\begin{dica}
**Sugestão de trabalho:** leia o capítulo na reunião de planejamento do 2º ano e adapte as rotas 1A–1D à turma.
\end{dica}

\chapter{A Técnica da Boquinha no 2º Ano}

No 2º ano, a Boquinha deixa de ser apenas visual e passa a ser metalinguística: o aluno articula, sente a posição dos articuladores, compara pares mínimos (R/L, S/Z, X/CH) e corrige a própria fala com apoio do espelho. Cada unidade traz a descrição articulatória dos sons trabalhados (foco, modo, vozeamento).


\begin{dica}
**Sugestão de trabalho:** leia o capítulo na reunião de planejamento do 2º ano e adapte as rotas 1A–1D à turma.
\end{dica}

\chapter{Alinhamento Curricular: BNCC e Matrizes Externas}

Este volume alinha-se às habilidades de Língua Portuguesa da BNCC para o 2º ano (EF02LP01--EF02LP10) e prepara o aluno para as matrizes de referência do SPAECE e dos itens CAEd (leitura, localização de informação, gêneros textuais, análise linguística e produção).


\begin{dica}
**Sugestão de trabalho:** leia o capítulo na reunião de planejamento do 2º ano e adapte as rotas 1A–1D à turma.
\end{dica}

\chapter{Níveis Internos de Progressão (K0–K5)}

A progressão interna organiza o avanço do escolar: da consolidação da base alfabética à fluência segmental, à ortografia regulada por regras e à fluência composicional. O aluno circula por rotas 1A–1D conforme o resultado dos checkpoints, sem rótulos nem diagnóstico.


\begin{dica}
**Sugestão de trabalho:** leia o capítulo na reunião de planejamento do 2º ano e adapte as rotas 1A–1D à turma.
\end{dica}

\chapter{Rotas Individualizadas de Aprendizagem (1A–1D)}

Rotas 1A (consolidação), 1B (avanço), 1C (desafio) e 1D (enriquecimento) orientam a diferenciação: escolha de planchetas, sequência de atividades de sons e ritmo das folhas Kumon, sempre a partir de evidências do Registro de Progresso.


\begin{dica}
**Sugestão de trabalho:** leia o capítulo na reunião de planejamento do 2º ano e adapte as rotas 1A–1D à turma.
\end{dica}

\chapter{Design Neuroinclusivo: Tipografia, Contraste e Ritmo}

Fonte cursiva escolar treinável, margens calmas, caixas amigáveis, instruções curtas e apoio visual. O design reduz sobrecarga cognitiva e favorece alunos com TDAH, dislexia e transtornos de linguagem, sem substituir a avaliação especializada.


\begin{dica}
**Sugestão de trabalho:** leia o capítulo na reunião de planejamento do 2º ano e adapte as rotas 1A–1D à turma.
\end{dica}

\chapter{Avaliação Formativa: Registro, Rastreio e Encaminhamento}

Três camadas: (a) registro de progresso diário; (b) checkpoints de rastreio entre unidades (autoavaliação do aluno e observação do professor, S/F/R/N); (c) avaliações diagnósticas do volume. Sinais persistentes são discutidos com a família e avaliados pelo profissional habilitado (Volume Profissional de Sondagem e Rastreio).


\begin{dica}
**Sugestão de trabalho:** leia o capítulo na reunião de planejamento do 2º ano e adapte as rotas 1A–1D à turma.
\end{dica}

\chapter{Preparação para SPAECE, CAEd e Avaliações Externas}

Itens-modelo no formato das matrizes externas, treino de comando ('Circule', 'Assinale', 'Localize') e de tempo de tarefa, com leitura mediada em sala e avaliações simuladas ao fim de cada trimestre.


\begin{dica}
**Sugestão de trabalho:** leia o capítulo na reunião de planejamento do 2º ano e adapte as rotas 1A–1D à turma.
\end{dica}

\chapter{Sequência Didática Granular: Lógica Kumon no 2º Ano}

Cada unidade fragmenta o objetivo em etapas microsequenciadas: 5 sessões por semana, 10 a 15 minutos por dia. A gradação é a alma do método: nenhum aluno avança sem consolidar a etapa anterior; o rastreio integrado garante a observação sistemática entre as etapas.


\begin{dica}
**Sugestão de trabalho:** leia o capítulo na reunião de planejamento do 2º ano e adapte as rotas 1A–1D à turma.
\end{dica}

\chapter{Ciência da Leitura Aplicada ao 2º Ano}

No 2º ano, a ciência da leitura (Dehaene, Salles, Capovilla) orienta a prática: leitura por decodificação automatizada, vocabulário, compreensão e fluência. O treino Kumon automatiza a via fonológica; a leitura compartilhada alimenta a compreensão; a escrita diária consolida a ortografia. O professor conhece a razão de cada etapa.


\begin{dica}
**Sugestão de trabalho:** leia o capítulo na reunião de planejamento do 2º ano e adapte as rotas 1A–1D à turma.
\end{dica}

\chapter{Fonologia do Português Brasileiro em Sala}

No 2º ano, o professor domina o mapa fonológico do PB: dígrafos (CH, LH, NH, RR, SS, QU, GU), encontros consonantais, nasalidade (AM/AN, ÃO, ÃE, ÕE) e polissemia do X. A Técnica da Boquinha traduz esse mapa em gestos articulatórios observáveis, e as listas de sons da Parte VII organizam o treino por fonema.


\begin{dica}
**Sugestão de trabalho:** leia o capítulo na reunião de planejamento do 2º ano e adapte as rotas 1A–1D à turma.
\end{dica}

\chapter{Planejamento Anual do 2º Ano}

O planejamento anual distribui as 21 unidades pelos quatro bimestres, intercala as folhas Kumon, prevê as avaliações diagnóstica, formativa e somativa e agenda os checkpoints de rastreio. O quadro de planejamento consta do Guia do Professor (Parte V); cada bimestre encerra com uma avaliação e uma reunião de pais.


\begin{dica}
**Sugestão de trabalho:** leia o capítulo na reunião de planejamento do 2º ano e adapte as rotas 1A–1D à turma.
\end{dica}

\chapter{Avaliação Formativa no Ciclo de Alfabetização}

A avaliação formativa acompanha o processo: observa a leitura em voz alta, o ditado da semana, a escrita espontânea e a folha Kumon do dia. O registro do professor (S/F/R/N) alimenta o Registro de Progresso e a reunião de pais. A avaliação somativa do fim do bimestre usa as atividades da Parte IV e o gabarito; a avaliação diagnóstica inicial e a triagem de risco usam as fichas do Volume Profissional. Nenhuma nota é usada para classificar: a nota é devolvida ao aluno em forma de meta de fluência e de escrita.


\begin{dica}
**Sugestão de trabalho:** leia o capítulo na reunião de planejamento do 2º ano e adapte as rotas 1A–1D à turma.
\end{dica}

\chapter{O Papel do Erro na Aprendizagem da Escrita}

O erro ortográfico revela a hipótese da criança sobre a escrita. Trocar P por B ou omitir o M antes de P não é 'falta de atenção': é uma regularidade a trabalhar com a Boquinha, a grade silábica e o ditado da semana. O erro deve ser corrigido com devolutiva imediata e repetição da mesma folha, nunca com punição. O professor que entende o padrão do erro planeja a próxima unidade com precisão; por isso cada checkpoint registra o padrão, não o aluno.


\begin{dica}
**Sugestão de trabalho:** leia o capítulo na reunião de planejamento do 2º ano e adapte as rotas 1A–1D à turma.
\end{dica}
